\section{Consequences and examples}\label{sec:consequences}

\subsection{Effective dissipation and logarithmic growth}

For the fixed inclusion define the effective cubic dissipation
coefficient
\[
 A(T)=\frac{-Z'(T)}{Z(T)^3}
     =\frac{2\mathcal E(T)}{Z(T)^3}.
\]
The adjective refers only to this scalar identity; it does not
assert that \(Z\) solves an autonomous scalar equation.

\begin{corollary}\label{cor:dissipation-coefficient}
Under the reduced, distinct-factor hypotheses,
\begin{equation}\label{eq:effective-coefficient}
 cT\le A(T)\le CT,\qquad
 \liminf_{T\to\infty}\frac{A(T)}T\ge C_*^2,\qquad
 \lim_{T\to\infty}\frac2{T^2}\int_0^T A(t)\,\dd t=C_*^2.
\end{equation}
On every compact subset of \(U\),
\begin{equation}\label{eq:logarithmic-rate}
 \log u(T)+\log T
 \longrightarrow \log\frac{v_0}{4\pi^2\sigma}
\end{equation}
uniformly. In particular \(\log Z(T)/\log T\to-1\).
\end{corollary}

\begin{proof}
The bounds follow from \(\mathcal E(T)\asymp T^{-2}\)
and \(Z(T)\sim C_*^{-1}T^{-1}\).
For the lower limit, write
\[
 \frac{A(T)}T
 =\frac{-Z'(T)}{Z(T)^2}\,\frac1{TZ(T)}
\]
and apply \eqref{eq:reduced-liminf}. The exact identity
\[
 Z(T)^{-2}=Z(0)^{-2}+2\int_0^T A(t)\,\dd t
\]
and the mass equivalent prove the averaged limit.
Finally \(u/Z\to1/v_0\) uniformly on each compact
subset of \(U\); its limit is positive. Combining
this with \eqref{eq:actual-mass} and taking logarithms
proves \eqref{eq:logarithmic-rate}.
\end{proof}

Neither \eqref{eq:effective-coefficient} nor the mass equivalent
asserts a pointwise equivalent for \(\mathcal E\) or \(A\).
The mass equivalent excludes a nonconstant leading logarithmic
multiplier in this class, but does not determine subleading terms.

There is a precise comparison with the two-weight description
in \cite[Section~3]{W_Singular}. In the category of equal-slope
semistable sheaves modulo sheaves supported at finitely many
points, the extension has subobject chain \(0<L<V\).
Its rank-polarized weight problem minimizes
\[
 w_L^2+w_M^2\quad\text{subject to}\quad w_M-w_L\ge1,
\]
giving \((-1/2,1/2)\), with no further iterated weights.
These are weights for squared metric norms. The proposed
line metric \(T^{-1/2}h_L\) consequently has the same
two-weight normalization as the finite model in
\cite[Example~5.12]{HKKP23}.
Formula~\eqref{eq:logarithmic-rate} gives the actual
exponent \(-1\) for the present singular class.
Thus the difference from \(-\frac12\log T\) is
\(-\frac12\log T+O(1)\), not merely an unspecified
unbounded negative remainder. The local ideal, its
Jacobian determinant in the fixed metrics, and the
matching through the compact extension supply information
which that polarized chain does not record.

\subsection{A toric tangent bundle}

We spell out the specialization to the tangent-bundle example
of \cite[Section~2]{W_Singular}. Let \(X\) have fan with cyclic rays
\[
 (1,0),\quad(1,1),\quad(0,1),\quad(-1,0),\quad(1,-1).
\]
Consecutive determinants are one. Removing \((1,1)\)
gives the first Hirzebruch surface, so \(X\) is its
blowup at the corresponding fixed point. With pullback
section and fibre classes \(S,f\) and exceptional class \(E\),
\[
 S^2=-1,\quad Sf=1,\quad f^2=0,\quad E^2=-1,\quad SE=fE=0.
\]
Choose a Kähler form in \(4S+7f-2E\), and put
\[
 V=T_X,\qquad L=\OO_X(2S+f-E),\qquad M=\OO_X(2f).
\]
The torus direction \((1,0)\) gives a saturated tangent
line. Its divisorial class is the sum of the two
opposite-ray divisors, namely \(2S+f-E\).
On the chart of the cone \(((1,1),(0,1))\), its
inclusion and quotient have the presentation
\[
 \iota=(z_1,-z_2)^{\mathsf t},\qquad q=(z_2,z_1).
\]
All other maximal cones contain one of the two
opposite rays; saturation there removes the divisorial
zero and leaves a nonvanishing generator. Thus the
only zero is the reduced fixed point \(p\), and
\[
 0\longrightarrow L\longrightarrow T_X
          \longrightarrow\cI_pM\longrightarrow0.
\]
The intersection pairing gives
\[
 \deg_\omega L=\deg_\omega M=8,\qquad
 v_0=18,\qquad \kappa=8\pi/9,\qquad
 c_2(T_X)-c_2(L\oplus M)=1.
\]
In particular the two factors are distinct and all
hypotheses of the positive theorem hold.
Writing \(\lambda=\lambda_p\) in the fixed normalization,
\begin{equation}\label{eq:toric-asymptotic}
 Z(T)\sim\frac{81\lambda}{\pi^2T},\qquad
 r_p(T)\sim\frac{3}{\sqrt2\,\pi}\,T^{-1/2}.
\end{equation}
These follow by substituting \(\sigma=\lambda^{-1}\).
The volume factor is evaluated here; the geometric
constant \(\lambda\) has not been numerically computed.
The radius formula is in base-orthonormal coordinates.
The stationary family obtained by changing the
polarization in \cite[Section~7]{W_Singular} is a
different family; its parameter is not flow time.

\subsection{Nonreduced points and repeated factors}

For arbitrary finite \(W\), Theorem~\ref{thm:general-upper-rate}
gives \(Z=O(T^{-1})\), and the presence of any nonreduced
point gives \(TZ\to0\). The proof supplies neither a
positive inner model nor a sharp nonreduced exponent.
The simple-zero normalization in
Lemma~\ref{lem:normalized-reference} is unavailable
there. An extension to that setting must retain
the full ideal germ, not just its length.

When \(L\simeq M\), the real small spectral space
has dimension three by \eqref{eq:weak-dimension}.
The projected correction and its parameter bounds
retain those directions. The scalar equation
\eqref{eq:comparison-ode}, and hence our positive
actual-flow comparison, are not proved for that case.
In particular, one cannot remove two directions
from the projection to infer the same scalar law.

Both exclusions occur in a concrete example. On a
projective K3 surface take a length-two curvilinear
scheme with local ideal \((z_1^2,z_2)\).
There is a locally free extension
\[
 0\longrightarrow\OO_X\longrightarrow V
                  \longrightarrow\cI_W\longrightarrow0.
\]
Indeed \(K_X=\OO_X\), and the Cayley--Bacharach
condition for the Serre construction holds: a constant
vanishing on a length-one subscheme vanishes on all
of \(W\). Apply \cite[Theorem~5.1.1]{HuybrechtsLehn97};
the local complete intersection condition is immediate
from the displayed regular sequence. The determinant
is trivial, \(c_2(V)=2\), and both factors have degree
zero for every Kähler form. The general nonreduced
conclusion applies, but the reduced-point profile
and radius theorem do not.

Finally, the positive result concerns a fixed compact surface
and a fixed holomorphic extension. It does not classify
singular graded sheaves of higher rank, varying holomorphic
defects, or singular fibres in a degenerating threefold family.
The centers here are fixed by \(W\); their leading radii
are linked through one amplitude and the constants
\(\lambda_p\). These are features of the stated problem,
not a classification of all possible concentration parameters.
